\section{Mechanism-specific force-kernel experiments}
These runs compare force-query implementations within the specified mechanism models. The continuous detection study in Section~\ref{sec:general-contact} uses the general STEP backend.

\subsection{Industrial mechanisms at matched accuracy}
\label{sec:industrial}
The three mechanism runs cover a driven cam and single follower, a 24/48-tooth spur pair, and a two-stage spur--helical reducer with 24/48 and 22/44 teeth. The gear input is driven at 1,800 r/min. The cam rotates at approximately 4 rad/s with elastic drive and a 3.2 N follower preload. Its moment of inertia is 0.007 kg\,m$^2$ and the follower has a mass of 0.48 kg. The cam uses the 12 mm activation layer and 5,000 N/m stiffness of the primitive cases. The cam runs for 20 s, the spur pair for 2 s and the reducer for 3 s. These durations cover approximately 13, 60 and 90 input revolutions, respectively. The gear law uses a 0.055 mm activation gap and a nominal stiffness of $1.6\times10^8$ N/m before axial and end weighting. The brake includes a load rise near 0.10 s, a release near 0.24 s, and a periodic component. The two-stage train uses one common shaft coordinate and two axial gear planes. The STEP models, material and load data are provided with the runs.

We compare the entire time histories against refined native solutions on common output grids of 1 ms for the cam and 20~\microsec{} for the gears. Position error is RMS follower displacement error or RMS transmission error at the driven base circle. Force error is the RMS force difference divided by the reference RMS force. The reducer must satisfy the joint tolerance in both stages. The native reference steps are 62.5~\microsec{} for the cam, 2.5~\microsec{} for the spur pair, and 5~\microsec{} with 63 helical sections for the reducer. The cam is refined from 1 ms and the gears from 20~\microsec{}. We evaluate three accuracy targets of 0.5~\micron{} and 2\%, 1~\micron{} and 5\%, and 2~\micron{} and 10\%. Table~\ref{tab:industrial} gives representative settings. The accompanying results retain the full ladders and sampled peak-force differences.

\begin{table}[pos=!htbp]
\centering\small
\caption{Representative mechanism settings. Position and force errors are RMS, and both error columns show the larger stage result for the reducer. Geometry settings refer to the number of cam polygon sides, the segments per gear section, or the mesher deflection. Online time excludes setup.}
\label{tab:industrial}
\begin{tabular}{llrrrrr}\toprule
Case & Geometry query & Step & Position & Force & Setup & Online \\
 & & (\microsec) & (\micron) & (\%) & (s) & (s) \\\midrule
Cam & Analytic cylinder--plane & 125 & 0.286 & 0.041 & 0.012 & 1.19 \\
 & Polygon support, 2048 & 125 & 0.175 & 0.024 & 0.018 & 1.59 \\
 & Triangle BVH, 2048 sides & 125 & 0.175 & 0.024 & 0.110 & 188.01 \\
 & OCCT distance per trial & 125 & 0.286 & 0.041 & 0.009 & 1108.26 \\
Spur & Native CAD & 5 & 0.037 & 0.89 & 0.017 & 4.03 \\
 & Polyline, 1024 & 5 & 0.037 & 0.89 & 0.018 & 3.94 \\
 & Triangle BVH, 391 nm & 5 & 0.256 & 3.35 & 3.832 & 47.81 \\
Reducer & Native CAD, 31 sections & 10 & 0.243 & 4.52 & 0.023 & 41.04 \\
 & Polyline, 256 & 10 & 0.244 & 4.55 & 0.028 & 41.45 \\
 & Triangle BVH, 1.563~\micron & 10 & 1.451 & 23.61 & 2.130 & 414.81 \\
 & Native CAD, 31 sections & 5 & 0.007 & 0.11 & 0.024 & 79.24 \\
\bottomrule\end{tabular}
\end{table}

\paragraph{Flat-follower cam.}
The rotating cylinder has a 105 mm radius and an eccentricity of 24 mm. The cam variants use the same HHT integrator, contact law and step guard. The analytic query computes the cylinder support point and checks its projection on the finite follower. The Open CASCADE variant calls \texttt{BRepExtrema\_DistShapeShape} on the original STEP shapes at each trial. It retains the reference shape's decomposition and uses a deflection setting of $10^{-7}$ mm. The two faceted variants use the same 2048-sided geometry. Polygon support selects the extremal side vertex and checks its projection on the follower, while the triangle BVH searches for a closest pair of points. The Open CASCADE variant makes 320,001 shape-distance calls across the 160,000 HHT steps. Its online time measures the entire integration. All four variants complete the 20 s run at 125~\microsec{} and meet the joint tolerance of 0.5~\micron{} and 2\%. The two faceted queries give the same RMS errors at the reported precision and differ substantially in online cost. Figure~\ref{fig:cam} compares the analytic and triangle BVH responses.

At the relaxed tolerance of 1~\micron{} and 5\%, a 512-sided polygon at 250~\microsec{} has 0.794~\micron{} displacement error and runs in 0.71 s in a separate batch. Partial cancellation between faceting and time-discretization error helps that setting meet the target. The refinement ladder measures this interaction across time steps and mesh resolutions.

The Drake comparison uses a cylinder with the 12 mm layer incorporated into its radius and a prismatic half-space. This preserves the maximum radial support height, force, eccentric torque, prescribed inertia and spring-damper loads of the model. At 125~\microsec{}, the Linux HHT and Drake RK2 runs meet the 0.5~\micron{} target with displacement errors of 0.286 and 0.402~\micron{}. Their median online times over five runs are 1.96 and 12.06 s, respectively. The point-level force law agrees within $1.5\times10^{-13}$ N along the Drake trajectories.

\begin{figure}[pos=!htbp]\centering
\includegraphics[width=.91\linewidth]{cam-response.pdf}
\caption{Eccentric-cam follower displacement and normal force from the analytic cylinder--plane calculation and the general triangle BVH at the same time step. The model is shown in Fig.~\ref{fig:sampled-cam}.}\label{fig:cam}\end{figure}

\paragraph{Single-stage spur pair.}
For the spur mesh we compare native flank sections against both a general triangle BVH and a faster piecewise-linear section representation. The section baseline preserves the same tooth-pair selection, sampled axial geometry, stiffness and line-of-action projection as the native method. At 5~\microsec{}, the native and 1024-segment polyline solutions meet the finest joint tolerance with nearly identical errors. The serial median times are 4.03 and 3.94 s and their repeat ranges overlap. The best observed triangle-BVH force error at this step is 3.35\% for the 391 nm mesh setting, which requires 47.81 s in the refinement batch. Figure~\ref{fig:spur} compares the native and tuned section solutions over the load-rise window. The force-waveform error at 10~\microsec{} already exceeds 5\% for the native method, so refining geometry alone does not achieve the tighter target.

\resultfigure{spur}{Single-stage spur pair and close-up of the tooth contact region. The transmission error and total mesh force show the load-rise interval at a 5~\microsec{} internal step. The two curves nearly coincide.}{fig:spur}

\paragraph{Two-stage reducer.}
The complex model places a spur stage upstream of a 22-degree helical stage. The native 31-section, 10~\microsec{} setting completes the full 3 s with worst-stage errors of 0.243~\micron{} and 4.52\%. The 256-segment section baseline at the same step gives 0.244~\micron{} and 4.55\%. In the serial batch, both methods take approximately 41 s and their repeat ranges overlap. The tested triangle settings resolve the helical stage more closely than the spur stage. At 5~\microsec{}, increasing the native axial quadrature from 31 to 63 sections changes the worst stage by 0.007~\micron{} and 0.11\%. The second stage is shown in Fig.~\ref{fig:reducer}. Direct engine-to-engine comparisons use the eccentric cam and the primitive contacts in Appendix~\ref{sec:analytic}, where a common native contact law can be checked from independent traces. The gearing comparisons hold the projected contact law fixed within the same HHT implementation.

\begin{figure}[pos=!htbp]\centering
\includegraphics[width=.91\linewidth]{reducer-response.pdf}
\caption{Second-stage responses of the two-stage reducer shown in Fig.~\ref{fig:sampled-reducer}. The window covers the applied load increase. The 31- and 63-section runs use the same 5~\microsec{} time step.}\label{fig:reducer}\end{figure}

\subsection{Effect of force enclosures and reuse}
\label{sec:ablation}
We compare the fixed native force kernel, the bounded kernel with tight evaluation, and the residual-driven mode with reuse. At 500~\microsec{} over 20 s, the cam's residual-driven mode with reuse records 113 omitted low-force pairs, 692 revalidated exclusions and 119 precision refinements in 40,000 accepted steps. In this batch the tight and adaptive cached bounded runs take 2.07 and 1.96 s. The original direct force computation takes 0.59 s. The three modes have a 4.479~\micron{} displacement RMS error relative to the 62.5~\microsec{} reference. Most of this error is associated with the integration step.

At a 10~\microsec{} step over 2 s, the spur gear's tight and cached bounded variants complete with the same reported trajectory and force errors as the fixed native kernel. The cached run obtains over $3.5\times 10^5$ accepted ray-coordinate enclosures from prior points but also reaches the tighter evaluation on most steps. Its online time is higher than that of the tight bounded mode in the paired ablation. In a separate 0.2 s diagnostic that reaches the first reducer load transient at a 20~\microsec{} internal step, both bounded variants complete all 10,000 steps. The cached version reuses over $9\times10^5$ ray-coordinate enclosures. It also requires 9,987 refinements and takes more time than the tight version in this diagnostic. These effects depend on how often a trial changes the active contacts or yields overlapping residual intervals.
\FloatBarrier
